\section{总结与延伸}

\subsection{一张表回收三套系统}

下面的表按候选来源、外部约束、中间 artifact 与主要 failure mode 对齐三项研究。这样比记住论文名更有用，因为后续任何“让大模型更可靠”的系统都可以用同一组问题审计。

\begin{longtable}{@{}L{0.13\textwidth}L{0.16\textwidth}L{0.19\textwidth}L{0.18\textwidth}L{0.19\textwidth}@{}}
\toprule
系统 & 候选来源 & 外部约束 & 可审计 artifact & 主要失败模式 \\
\midrule
Maieutic & LM 正反 explanations & 逻辑 clauses、Weighted MaxSAT & explanation graph、assignment & 候选知识错误、搜索成本、解释不 faithful \\
Symbolic KD & GPT-3 knowledge generations & RoBERTa critic、threshold、relation schema & filtered graph、student corpus & critic bias、稀有知识丢失、teacher cost \\
Delphi & norm data 与 commonsense backbone & task label、data policy、免责声明 & judgment、dataset provenance & 文化偏差、adversarial prompt、authority misuse \\
Delphi Hybrid & reasons、retrieved rules、norm scores & top-down principles、graph consistency & argument graph、audit trace & rule conflict、retrieval error、latency \\
\bottomrule
\end{longtable}

\subsection{最终结论}

第一，语言流畅性、知识质量、逻辑一致性与伦理判断必须分开验收。第二，结构化中间层的价值不仅是提高 accuracy，还包括让数据、理由、约束和失败可见。第三，“smaller and better”通常来自任务专门化、筛选与接口设计，不是参数减少本身。第四，machine ethics 的最高风险不是单个错误答案，而是把 prediction 偷换成 authority。第五，价值多元与制度责任不能被一个 aggregate label 取代。

\begin{importantbox}{最值得带走的一句话}
大语言模型可以是候选知识与候选理由的强大生成器，但可靠系统必须另外回答：谁来筛选、依据什么约束、保存什么证据、怎样拒答，以及谁拥有最终决策权。
\end{importantbox}

\subsection{拓展阅读}

建议按课堂依赖顺序阅读。先读 Maieutic Prompting，理解 explanation graph 与 consistency；再读 Symbolic Knowledge Distillation，理解 critic 如何把大模型输出变成 corpus；最后读 Delphi，重点阅读 data、limitations、social justice 与 counterarguments，而不是只看主结果表。

\begin{itemize}
  \item Jung et al., \href{https://arxiv.org/abs/2205.11822}{Maieutic Prompting: Logically Consistent Reasoning with Recursive Explanations}。
  \item West et al., \href{https://arxiv.org/abs/2110.07178}{Symbolic Knowledge Distillation: from General Language Models to Commonsense Models}。
  \item Jiang et al., \href{https://arxiv.org/abs/2110.07574}{Can Machines Learn Morality? The Delphi Experiment}。
  \item Bosselut et al., \href{https://aclanthology.org/P19-1470/}{COMET: Commonsense Transformers for Automatic Knowledge Graph Construction}。
  \item Hwang et al., \href{https://arxiv.org/abs/2010.05953}{COMET-ATOMIC 2020}。
\end{itemize}

\subsection{课后自检}

\begin{importantbox}{十二个检查问题}
\begin{enumerate}
  \item 为什么一个 Winograd 题答对不能证明 commonsense 已解决？
  \item Maieutic tree 为什么同时展开命题与否定命题？
  \item Weighted MaxSAT 改善 consistency，却为什么不能保证 truth？
  \item dataset solving、task solving 与 capability 有何差别？
  \item ATOMIC 的 relation types 如何改变数据收集和评估？
  \item COMET 与 GPT-3 的比较为什么不是 apple-to-apple？
  \item Symbolic KD 中 teacher、critic、corpus、student 各自承担什么资源成本？
  \item critical teacher 为什么可能让 student 超过 loose teacher？
  \item Delphi 的 descriptive judgment 与 normative ethics 有何区别？
  \item COMMONSENSE NORM BANK 的规模为什么不能证明文化中立？
  \item neuro-symbolic hybrid 增加了哪些可审计接口，也增加了哪些失败点？
  \item 为什么 value pluralism 需要 humanities 和制度参与，而不只是更大模型？
\end{enumerate}
\end{importantbox}

\end{document}
